\chapter{Hip-Fracture Surgery}

\section*{Introduction \& Clinical Indications}
Hip fractures include femoral-neck, intertrochanteric, and subtrochanteric injuries. In older adults they commonly result from a fall in fragile bone and carry substantial risks from immobility, delirium, thrombosis, pneumonia, pressure injury, and loss of independence. Surgery is usually performed after rapid medical optimization, with the aim of restoring stability and enabling early mobilization.

Displaced femoral-neck fractures are often treated with hemiarthroplasty or total hip arthroplasty. Stable intertrochanteric fractures may be treated with a sliding hip screw or cephalomedullary device, while unstable or subtrochanteric patterns generally require a cephalomedullary nail or another construct. Treatment depends on fracture pattern, pre-injury function, bone quality, cognition, comorbidities, and goals of care.

\section*{Relevant Surgical Anatomy}
The femoral head receives much of its blood supply through vessels traveling along the femoral neck, so displaced intracapsular fractures risk avascular necrosis and nonunion. The greater and lesser trochanters anchor important abductor and iliopsoas muscles. The sciatic nerve, femoral neurovascular bundle, profunda vessels, and hip capsule must be protected during positioning, reduction, and fixation.

\section*{Preoperative Workup \& Preparation}
Confirm the fracture with radiographs and use CT when the pattern is unclear or associated injuries are suspected. Assess pain, skin, neurovascular status, anemia, anticoagulation, renal function, infection, cognition, nutrition, and cardiopulmonary risk. Provide analgesia, pressure-area protection, delirium prevention, and venous-thromboembolism planning. Explain bleeding, infection, fixation failure, nonunion, avascular necrosis, dislocation, leg-length difference, reoperation, and mortality risk.

\section*{Surgical Technique \& Procedure}
After anesthesia and careful positioning, the fracture is reduced and stabilized with the selected implant, or the femoral head is replaced when arthroplasty is indicated. Fluoroscopy confirms alignment, screw or nail position, and implant length. Cemented or uncemented stems are selected according to patient and surgical factors. Stable fixation and a plan for weight-bearing are documented before closure.

\section*{Complications \& Risks}
Early complications include delirium, chest infection, urinary infection, venous thromboembolism, bleeding, wound infection, nerve or vessel injury, and perioperative cardiac events. Later complications include nonunion, avascular necrosis, cut-out or peri-implant fracture, malunion, leg-length inequality, dislocation, loosening, and need for revision. Recurrent falls and osteoporosis must be addressed.

\section*{Postoperative Care \& Recovery}
Mobilize early with physiotherapy and occupational therapy, using weight-bearing instructions appropriate to the fixation. Provide thromboprophylaxis, multimodal analgesia, nutrition, pressure-area care, delirium prevention, and osteoporosis assessment. Review wound and hemoglobin, and coordinate discharge, home safety, walking aids, and follow-up imaging.

\section*{Outcomes \& Prognosis}
Earlier surgery after medical optimization is associated with better recovery for many patients, but urgent anesthesia must still be individualized. Functional outcome depends on pre-fracture mobility, cognition, rehabilitation, fracture type, complications, and social support. A coordinated orthogeriatric pathway improves comprehensive care.

\section*{References}
\begin{enumerate}
\item American Academy of Orthopaedic Surgeons. Management of Hip Fractures in Older Adults Clinical Practice Guideline.
\item National Institute for Health and Care Excellence. Hip fracture management.
\item British Orthopaedic Association. Standards for Trauma and Orthopaedics.
\end{enumerate}
